\section*{Chapter 5: Neutron Stars: Wonders Of Physics}
\begin{description}[style=multiline,leftmargin=2em,labelwidth=1.5em]

% DOCX paragraph 2039
\item[1.]\hypertarget{note:05-01}{}Two fundamental quantities characterise a wave phenomenon, such as that of
electromagnetic waves: the wavelength (i.e., the distance between two ``peaks'' in the
wave) and the frequency (i.e., the number of peaks arriving in a given interval of
time). The wavelength, \emph{λ}, and the frequency, \emph{f}, of the electromagnetic
radiation\index{radiation!electromagnetic} are inversely proportional. More precisely, they are
related as \emph{λ = c / f}, where \emph{c} is the speed of light.

% DOCX paragraph 2040
\item[2.]\hypertarget{note:05-02}{}Radiation of this type is also called \emph{soft X-ray}\index{X-ray!radiation},
distinguishing it from radiation with higher energy which is defined as \emph{hard
X-ray}\index{X-ray!radiation}. The hardness of a photon in the X-ray\index{X-ray!radiation} band
is therefore measured through its energy (or, equivalently, its frequency). Soft X-rays
usually have energy between 0.1 and 0.3 keV (kiloelectronvolt), while hard X-rays
between 10 and 100 keV.

% DOCX paragraph 2041
\item[3.]\hypertarget{note:05-03}{}I should recall that first classifications grouped the types of stars into
\emph{``colour classes''}, that is, according to the dominant colour in which their
emission spectrum\index{spectrum!electromagnetic} appeared. Hence, stars were ordered from the
blue-white of the ``hottest'' stars to the orange-red of the ``coldest''. In a
decreasing order of temperature, or, equivalently from blue to red, the classes were: O,
B, A, F, G, K, M. To help to remember the sequence there is a simple phrase introduced
by the American astronomer who worked on this classification, Annie Jump Cannon
(1863-1941): \emph{``Oh, be a fine girl/guy, kiss me!''}.

% DOCX paragraph 2042
\item[4.]\hypertarget{note:05-04}{}To fully understand why it was a bizarre phenomenon, it would be good to deepen the
question a little. It is natural to expect that the timescale of
variability\index{variability} of a celestial source is proportional to its size. In
particular, the brightness of larger objects can only vary slowly, while that produced
by smaller objects (albeit on an astronomical scale!) can change more rapidly. This
relation between size and timescale is the consequence of a fundamental requirement: the
source is causally connected. In other words, its changes cannot be arbitrarily rapid
but can vary on a timescale (\emph{∆T} ) which must be greater than the time a photon
moving at the speed of light (\emph{c}) would take to travel through a source of a given
size. For this reason, this timescale is called the ``light-crossing'' time. At a
mathematical level, this constraint translates into a precise formula: if \emph{R} is
the characteristic size of the source, then to preserve causality we must impose that
\emph{R \textless{} c ΔT}. Since \emph{∆T} is measured by the observations, we easily
obtain an upper bound for the size of the source \emph{R}. To give you some reference,
the light-crossing time of the Sun is about five seconds, while its brightness is
subject to variations that occur on a much longer timescale of about 11 years. Hence,
the Sun must be smaller than 11 light years, which is obviously correct!

% DOCX paragraph 2043
\item[5.]\hypertarget{note:05-05}{}I. Shklovsky, \emph{Soviet Astronomy}, 11, 749 (1968)

% DOCX paragraph 2044
\item[6.]\hypertarget{note:05-06}{}W. Baade, F. Zwicky, \emph{Proceedings of the National Academy of Sciences of the
United States of America}, 20, 254 (1934)

% DOCX paragraph 2045
\item[7.]\hypertarget{note:05-07}{}A. Hewish\index{Hewish, Antony}, S. J. Bell, J. D. H. Pilkington, P. F. Scott, R. A.
Collins, \emph{Nature}, 217, 709 (1968)

% DOCX paragraph 2046
\item[8.]\hypertarget{note:05-08}{}The \emph{atomic mass} is the number that indicates how many neutrons and protons are
contained within a certain atomic nucleus. A similar but distinct quantity is the
\emph{atomic number}, which instead only counts the number of protons in an atomic
nucleus. For example, hydrogen has atomic mass 1 and atomic number 1 because its nucleus
is made up of one proton only; helium has atomic mass 4 and atomic number 2 since its
nucleus is made up of two protons and two neutrons.

% DOCX paragraph 2047
\item[9.]\hypertarget{note:05-09}{}Neutrons are fermionic particles, i.e., they follow a Fermi--Dirac statistical
distribution and are therefore subject to the Pauli exclusion principle. This principle
establishes that a quantum state can be occupied at most by a single particle of this
type. This very principle is responsible for the degeneracy pressure\index{pressure!degeneracy}.
Using an analogy that has unfortunately become all too clear these days, it is as if
fermionic particles were forced to ``socially distance'' and therefore could not be too
close to each other. When they do get too close, they generate a lot of pressure that
prevents further compression.

% DOCX paragraph 2048
\item[10.]\hypertarget{note:05-10}{}In quantum mechanics\index{quantum mechanics}, the wave function represents the state of a
physical system of, say, an elementary particle. It is a function of time and position,
so it can be used to measure the probability that a particle is in a certain region of
space at a given time.

% DOCX paragraph 2049
\item[11.]\hypertarget{note:05-11}{}G. Gamow, M. Schoenberg, \emph{Physical Review}, 59, 539 (1941)

% DOCX paragraph 2050
\item[12.]\hypertarget{note:05-12}{}The model I proposed is called ``\emph{blitzar}'' and involves the birth of a black
hole from a neutron star that collapses, losing its magnetic field. Numerical
simulations of this process show that it would produce a signal very close to that
observed, both in duration and in the amount of energy emitted. Obviously, this does not
mean that the model is correct or that it really explains the observations: the jury is
still out~\ldots{} (see H. Falcke and L. Rezzolla, in \emph{Astronomy and Astrophysics},
562, 134 (2014)).

% DOCX paragraph 2051
\item[13.]\hypertarget{note:05-13}{}F. Pacini, in \emph{Nature}, 219, 145 (1968)

% DOCX paragraph 2052
\item[14.]\hypertarget{note:05-14}{}I have decided to provide an example here that is different from PSR J0437-4715 to
clarify that these are not ``special'' pulsars.

\end{description}
